\section{第16组：固体能带与激光}
\sourceproblem{16-T1}{104}{108}
\problemtext 从能带结构说明绝缘体与半导体导电性能差异。
\workings 理想本征材料在零温时价带满、导带空。半导体禁带较窄，有限温度热激发可产生导带电子与价带空穴，载流子浓度有 $\exp[-E_g/(2k_BT)]$ 的主要温度依赖；绝缘体通常禁带较宽，同温度下热激发载流子很少。
\crosscheck 能隙并非唯一因素，迁移率、杂质、缺陷、温度均影响实际电导；没有对所有材料都适用的严格能隙分界值。
\sourceproblem{16-T2}{104}{108}
\problemtext 本征半导体锗禁带宽0.67 eV，估算太阳电池带间吸收的最大波长。
\workings 光子能量需满足 $hc/\lambda\ge E_g$，故 $\lambda_{\max}=hc/E_g\approx1.85$ 微米。
\crosscheck 是能量阈值估计；锗为间接带隙材料，实际带边吸收还涉及声子和跃迁概率。
\sourceproblem{16-T3}{105}{109}
\problemtext 解释粒子数反转及实现的必要条件。
\workings 对激光跃迁上、下能级，反转条件为 $N_u/g_u>N_l/g_l$（$g$ 为简并度），等简并时为 $N_u>N_l$。正常正温热平衡不满足该条件，须由外界泵浦建立非平衡分布，并使抽运及弛豫机制能够维持反转。
\crosscheck 粒子数反转只给受激净增益的必要基础；持续激光振荡还要求增益克服损耗。
\sourceproblem{16-T4}{105}{109}
\problemtext 说明激光器的主要组成和光学谐振腔作用。
\workings 通常包括增益介质、泵浦系统和光学反馈谐振腔。谐振腔提供多次通过增益介质的反馈，选择满足相位和损耗条件的模式，并经部分透射镜输出光。稳态振荡须满足往返净增益补偿损耗。
\crosscheck 泵浦供能，谐振腔本身不产生能量；无腔放大器也能有受激辐射，但不等同于通常的谐振腔激光振荡器。
